\documentclass[11pt,letterpaper]{article}

\usepackage[english]{babel}
\usepackage[T1]{fontenc}
\usepackage{lmodern}
\usepackage[margin=2.6cm,top=2.3cm,bottom=2.4cm]{geometry}
\usepackage{microtype}
\usepackage{amsmath}
\usepackage{booktabs,tabularx,array}
\usepackage{needspace}
\usepackage[section]{placeins}
\usepackage{listings}
\usepackage{xcolor}
\usepackage{xurl}
\usepackage{fancyhdr}
\usepackage[colorlinks=true,linkcolor=blue,citecolor=blue,urlcolor=blue]{hyperref}

\hypersetup{pdftitle={ffmt: Freestanding-Friendly C++ Formatting and Printing for GPUs},
  pdfsubject={Design, portability constraints, and a validation plan}}
\lstset{language=C++,basicstyle=\small\ttfamily,
  keywordstyle=\bfseries,commentstyle=\itshape,
  showstringspaces=false,columns=fullflexible,keepspaces=true,
  breaklines=true,frame=single,rulecolor=\color{black!20},
  backgroundcolor=\color{black!2},framerule=0.3pt,
  xleftmargin=0.5em,xrightmargin=0.5em,
  aboveskip=0.8em,belowskip=0.8em}
\newcolumntype{Y}{>{\raggedright\arraybackslash}X}
\newcolumntype{P}[1]{>{\raggedright\arraybackslash}p{#1}}
\renewcommand{\arraystretch}{1.1}
\newcommand{\ffmt}{\texttt{ffmt}}
\newcommand{\fmtlib}{\texttt{\{fmt\}}}
\setlength{\parskip}{0.25em}
\setlength{\parindent}{0pt}
\setlength{\headheight}{14pt}
\pagestyle{fancy}
\fancyhf{}
\fancyhead[L]{\small\ffmt: GPU-friendly formatting}
\fancyhead[R]{\small Working draft / Revision 2}
\fancyfoot[C]{\thepage}
\renewcommand{\headrulewidth}{0.3pt}
\emergencystretch=2em
\renewcommand{\topfraction}{0.90}
\renewcommand{\bottomfraction}{0.80}
\renewcommand{\textfraction}{0.08}
\renewcommand{\floatpagefraction}{0.80}


\title{\texttt{ffmt}: Freestanding-Friendly C++\\Formatting and Printing for GPUs}
\author{Me}
\date{21 October 2026}


\newcommand{\code}[1]{\texttt{#1}}

\begin{document}
\maketitle
\thispagestyle{plain}

\begin{abstract}
Modern C++ formatting provides readable, type-directed diagnostics, but its use
inside accelerator kernels depends on both character conversion and device
output facilities. We present the design of \ffmt, a C++20 implementation
that separates a brace-formatting interface from backend-specific emission.
A buffer path converts values into bounded temporary storage, while a
specifier path translates supported fields into a device \texttt{printf}
format and forwards runtime arguments. This separation exposes a portability
trade-off: controlling the generated bytes requires a usable string-output
primitive, whereas delegating conversion inherits runtime formatting and
address-space restrictions. We explain the compile-time interface, reference
semantics, SPIR-V constraints, floating-point considerations, and storage costs.
We organize preliminary author-reported testing into a qualification matrix
and define reproducible correctness, negative-compilation, and performance
experiments across host C++, SYCL, OpenMP target, CUDA C++, HIP C++, and Kokkos
execution spaces. We distinguish inspected implementation features from
unverified configurations and proposed extensions. The contribution is a
portability design and evaluation methodology; this working draft does not
claim complete standard-format conformance, strict freestanding conformance,
or measured performance superiority.
\end{abstract}

\section{Motivation and contribution}

C++20 introduced \texttt{std::format}, and C++23 added \texttt{std::print}
\cite{format-intro,print-intro}. Formatting constructs characters; printing
also delivers them to an output destination. GPU diagnostics need both, but
they need not use the same implementation or execute conversion in the same
place. A familiar format syntax alone does not establish device-callable
code, a usable address space, or reliable output delivery.

The goal of \ffmt\ is to let programmers write brace-style diagnostics inside
kernels without manually maintaining a parallel list of percent directives.
Its architecture is organized around three obligations: validate the format
and argument types, produce the intended representation, and deliver that
representation through the selected backend. Separating these obligations
makes failures attributable to the library, compiler, or runtime rather than
to an undifferentiated notion of ``GPU support.''

This article describes a two-path implementation and identifies the experiments
needed to qualify its semantics and portability. It is not a proposal for a
new decimal conversion algorithm. Related foundations include \fmtlib,
Dragonbox, the standard formatting facilities, and existing device-output
interfaces \cite{ffmt-readme,dragonbox,kokkos-printf,cuda-printf}.

\Needspace{6\baselineskip}
\paragraph{Evidence boundary.}
Implementation observations refer to the publicly accessible README, main
header, and Makefile inspected on 21 October 2026. A complete checkout was not
obtained, and no \ffmt\ test suite or accelerator execution was independently
run for this revision. The author's previously reported executions remain
separate from source inspection and proposed experiments. Independent host
reference probes accompany this revision, but are not \ffmt\ validation.

\section{Scope and compile-time interface}

\subsection{Freestanding-friendly, not a conformance claim}

Freestanding is a defined C++ implementation category, not a synonym for
header-only code, absence of allocation, or successful compilation in one
kernel \cite{compliance}. Here, freestanding-friendly means designing a core
whose storage, error handling, and output dependencies can be controlled.
Strict qualification would require a specified language/library baseline,
transitive dependency and symbol audits, and a minimal sink implementation.
Device runtime I/O is still a dependency even when the formatter itself avoids
heap allocation and exception unwinding.

It is inaccurate to explain every \fmtlib\ porting obstacle by allocation or
exceptions: output-iterator formatting and exception-disabled configurations
already exist \cite{fmt-api}. The relevant questions are which functions are
device callable, which headers and runtime symbols are needed, how storage is
provided, and how the output runtime consumes pointers and values.

\subsection{Known format text; runtime values}

The intended user interface is illustrated below. These examples describe the
API; they are not new device test results.

\begin{lstlisting}
#include <ffmt/base.hpp>

// Within a supported host or kernel context:
int b = 7;
FFMT_PRINTLN("Hello {} {} {:23.2f}", "world", b, 12.12);

int x = 7;
float y = 1.25f;
FFMT_PRINTLN("Hello {1:.2f} {0:d}", x, y);
\end{lstlisting}

The format text is captured during compilation; \texttt{x}, \texttt{y}, and
\texttt{b} may be runtime values. In particular, \texttt{\{:23.2f\}} specifies
width 23, precision 2, and fixed presentation; \texttt{\{f:23\}} is not its
positional-format equivalent. This project constraint should not be confused
with a universal limitation of standard formatting: \texttt{std::vformat}
accepts runtime format text \cite{format-functions}.

The inspected macro selects ordinary-argument syntax on the buffer path and
a string non-type template argument on the specifier path \cite{ffmt-source}:

\begin{lstlisting}
// Illustrative expansions of FFMT_PRINT, not new overloads:
ffmt::print("Hello {1:.2f} {0:d}", x, y);    // buffer path
ffmt::print<"Hello {1:.2f} {0:d}">(x, y);    // specifier path
\end{lstlisting}

A non-type template argument makes the characters available while generating
static output text, selecting argument indices, and instantiating backend
checks. Ordinary function parameters do not automatically become template
arguments merely because a caller supplied a literal. The macro hides this
implementation difference. It does not extend namespace \texttt{std};
\texttt{std::print<"...">} is not the API being implemented here.
A type-encoding user-defined literal could provide another interface, but
\texttt{"..."\_fmt} is a proposal, not an advertised existing overload.
Repeated positional arguments must reuse captured values without evaluating
an argument expression twice; this does not impose left-to-right evaluation
beyond the C++ language rules.

\subsection{Why not simply reuse \texttt{std::basic\_format\_string}?}

\texttt{std::basic\_format\_string} provides compile-time validation for its
argument types and access to the underlying string view; it does not expose
a standardized parsed instruction sequence, a generated \texttt{printf}
format, or an address-space-qualified storage object \cite{format-string}.
Its validation also follows \texttt{std::formatter} rather than automatically
understanding \ffmt's custom formatter protocol or backend restrictions.

Therefore reuse is worth evaluating as an optional host-side checker or API
adapter, but would not eliminate \ffmt's lowering and device-compatibility
work. A device-facing dependency on \texttt{<format>} also needs independent
toolchain qualification. This is an architectural rationale, not a verified
account of the original author's decision. Compile-time checking and
compile-time specialization are separate properties: neither alone proves
that all parsing metadata disappears from the generated kernel.

\section{Two paths and the output boundary}

The buffer path converts values into a fixed-capacity character buffer. The
specifier path generates a supported percent-format string and delegates
conversion to a device-output facility \cite{ffmt-readme}. The difference is
semantic, not merely a choice of optimization.

\begin{table}[!htbp]
\centering\small
\begin{tabularx}{\linewidth}{@{}P{2.65cm} Y Y@{}}
\toprule
Question & Buffer path & Specifier path \\
\midrule
Conversion owner & Library formatting code & Output runtime \\
Output-boundary payload & Produced bytes & Generated format and typed values \\
Primary benefit & Control over representation & Works with a narrower numeric-output interface \\
Primary obligations & Bound every write; transport bytes faithfully & Validate grammar, casts, ordering, and runtime semantics \\
Primary costs & Temporary storage and conversion code & Runtime dependence and a restricted supported subset \\
\bottomrule
\end{tabularx}
\caption{Architectural responsibilities, not measured performance results.}
\label{tab:paths}
\end{table}

For the reordered example, a compatible specifier lowering is conceptually:
\begin{lstlisting}
// Conceptual C-printf equivalent; actual runtime ABI may differ.
printf("Hello %.2f %d\n", static_cast<double>(y), x);
\end{lstlisting}
The directive sequence follows output order, not source argument order.
The lowering must also escape literal percent signs, select integer length
modifiers and signedness correctly, and handle repeated fields. Default
boolean text is not a portable device \texttt{\%b} directive. Likewise,
\texttt{\%b} should not be assumed available for binary integers simply because
a recent host C library supports it.

A byte-preserving sink might use \texttt{printf("\%s", buffer)} or
\texttt{fputs(buffer, stdout)}. In those calls percent characters in the data
are not parsed as directives. The inspected AdaptiveCpp hook is different:
its device route can pass the completed buffer as the format of an internal
print primitive and therefore performs percent escaping \cite{ffmt-source}.
A common front end does not imply an identical sink contract. Tests must cover
escaping expansion and truncation together, including a percent character at
the capacity boundary; an incomplete escaped directive is not ordinary
text truncation.

\subsection{Programming model, lowering, and hardware are different axes}

Table~\ref{tab:primitives} identifies the relevant layers. It does not claim
that every route has been qualified for \ffmt. OpenMP target is not a new
percent-format specification; the selected compiler and offload runtime
supply its output implementation. HIP C++ is an established programming
interface, not a speculative name \cite{hip}.

\begin{table}[!htbp]
\centering\small
\begin{tabularx}{\linewidth}{@{}P{2.7cm} Y Y@{}}
\toprule
Route & Output mechanism / lowering & Qualification concern \\
\midrule
Host C++ & Host C library sink & libc, locale, ABI, and architecture \\
Intel SYCL / SPIR-V & Intel experimental \texttt{printf}; OpenCL-style operation & Address spaces, compiler transformations, runtime support \\
CUDA C++ / NVPTX & Device \texttt{printf}, typically a \texttt{vprintf} runtime boundary & Device transport plus host-side final formatting \\
HIP C++ / AMDGPU & Compiler lowering to device-library buffered output or hostcall & Compiler version, selected transport, string copying \\
OpenMP target & Target-specific device runtime & Record SPIR-V, NVPTX, or AMDGPU separately \\
Kokkos & \texttt{Kokkos::printf} and selected execution space & Framework wrapper plus underlying backend, not a new ISA \\
\bottomrule
\end{tabularx}
\caption{Output layers. Sources: Intel, NVIDIA, LLVM, and Kokkos documentation
\cite{intel-printf,cuda-printf,amdgpu-printf,kokkos-printf}.}
\label{tab:primitives}
\end{table}

The correct Intel extension name is
\texttt{sycl::ext::oneapi::experimental::printf}, not
\texttt{sycl::external::printf}. Intel's 2026 documentation describes removal
of the old C-varargs implementation in favor of a variadic template that
preserves argument types \cite{intel-printf}. C ellipsis promotes
\texttt{float} to \texttt{double}; a variadic template need not. The inspected
\ffmt\ cast helper itself converts floating arguments to \texttt{double}
\cite{ffmt-source}. Thus float-only input must not be advertised as requiring
no device FP64 capability without an end-to-end check and, potentially, a
backend-specific cast policy.

LLVM's AMDGPU lowering contains both buffered and hostcall machinery, including
handling for nonconstant format text \cite{amdgpu-printf}. This is evidence
about that implementation, not a guarantee for every HIP version or target.
For \ffmt, dynamic \emph{data} passed to a fixed \texttt{\%s} is the priority;
dynamically generated format text is a separate capability. chipStar, which
runs HIP/CUDA-style programs through SPIR-V, is another distinct route
\cite{chipstar}; a chipStar result does not establish native AMD HIP support.

Grace Hopper pairs an Arm-based Grace CPU with a Hopper GPU
\cite{grace-hopper}. Record both sides and their toolchains in the matrix.
``DGX'' is a system designation, not a fourth device-output primitive;
a Grace-hosted GPU run also does not establish native Arm CPU formatting
behavior.

\section{Reference semantics and compatibility tables}

\subsection{A versioned reference, not an unconditional oracle}

The repository states a preference for \texttt{std::format} when it differs
from \texttt{fmt::format} \cite{ffmt-readme}. The semantic target should
therefore be a named standard wording baseline plus documented extensions,
with \fmtlib\ and standard-library implementations used as independent
comparison tools. Fixing a version matters: the current draft's default
floating-point notation rule differs from the historical C++23 draft rule
\cite{current-charconv,p3505}.

\begin{table}[!htbp]
\centering\small
\begin{tabularx}{\linewidth}{@{}P{2.8cm} Y Y@{}}
\toprule
Comparison & Standard baseline & \fmtlib\ / implication \\
\midrule
Default \texttt{double(1e15)} & N4950: \texttt{1e+15} & Historical exponent-threshold behavior: \texttt{1000000000000000}. Version-sensitive, not a timeless difference. \\
Named replacement fields & No named-argument facility in the cited standard grammar & \texttt{fmt::arg} is a library extension; not required for standard-subset agreement. \\
Compile-time checking versus compilation & Checked format-string construction does not mandate a particular specialized lowering & \texttt{FMT\_COMPILE} explicitly compiles format text; treat it as a separate benchmark configuration. \\
Runtime format text & \texttt{std::vformat} provides a runtime route & Runtime APIs also exist in \fmtlib; \ffmt's literal restriction is its own design choice. \\
\bottomrule
\end{tabularx}
\caption{Semantic/API distinctions, not a newly executed \fmtlib\ comparison.
Sources: standard wording and \fmtlib\ documentation
\cite{n4950-format,n4950-charconv,p3505,fmt-api,format-functions}.}
\label{tab:reference}
\end{table}

The \texttt{1e15} case is useful only with explicit version labels. Keep
expected bytes, argument types and bit patterns, reference versions, and the
intended wording alongside every regression. A disagreement may be a wording
change, intentional extension, implementation defect, or backend limitation.
Agreement with one library is not a conformance proof.

\subsection{What delegation to \texttt{printf} changes}

Table~\ref{tab:delegation} separates cases that translate naturally from those
that require special treatment or rejection. Features unsupported by the
current translator must not be described as fundamentally impossible in
all device runtimes.

\begin{table}[!htbp]
\centering\small
\begin{tabularx}{\linewidth}{@{}P{3.0cm} Y Y@{}}
\toprule
Case & Brace-format expectation & Specifier-path issue / policy \\
\midrule
\texttt{\{\}}, \texttt{1.23456789} & \texttt{1.23456789} & Plain \texttt{\%g} gives \texttt{1.23457}; not shortest-default equivalence. \\
\texttt{\{:.2f\}}, \texttt{1.25} & \texttt{1.25} & \texttt{\%.2f} agrees in the host probe; qualify other values and runtimes. \\
\texttt{\{:x\}}, \texttt{-42} & \texttt{-2a} & Unsigned \texttt{\%x} does not preserve the sign-and-magnitude spelling; reject or implement it. \\
\texttt{\{:a\}}, \texttt{1.5} & \texttt{1.8p+0} & Host \texttt{\%a} includes \texttt{0x}; current non-buffer contract rejects this presentation. \\
\texttt{\{:b\}}, \texttt{42} & \texttt{101010} & No assumed common device directive; currently buffer-only. \\
Centering / custom fill & Requested layout & Currently buffer-only; simple percent flags are insufficient. \\
Dynamic width / precision & Runtime integer controls layout & Current non-buffer contract rejects it; \texttt{*} support alone does not implement all brace semantics. \\
Runtime-produced string & Bytes from the supplied string & Depends on the sink; literal-only \texttt{\%s} cannot consume an arbitrary private buffer. \\
\bottomrule
\end{tabularx}
\caption{Delegation hazards and current documented exclusions
\cite{ffmt-readme,n4950-format,opencl-c}. Numeric spellings were independently
checked with GCC 14.2.0 / libstdc++ date 20250315 in the C locale; no \ffmt\ or
\fmtlib\ execution is implied.}
\label{tab:delegation}
\end{table}

The default-float case needs a contract decision: the README describes a
restriction, while the inspected implementation can map an unspecified
floating presentation to \texttt{g} \cite{ffmt-source}. A reduced-precision
fallback should be explicit, or the call should be rejected; it must not be
counted as exact semantic agreement.

The original trailing-zero example also needs precision: ordinary default
formatting of \texttt{2.0} produces \texttt{2}, whereas \texttt{\%f} produces
\texttt{2.000000} in the host probe. Both \texttt{\{:\#.0f\}} and
\texttt{\%\#.0f} produce \texttt{2.} there. The distinction is presentation
and precision, not a general rule that \fmtlib\ always removes zeros and
\texttt{printf} never does.

\section{SPIR-V: string capabilities and compiler behavior}

Three properties must be distinguished: C++ \texttt{const} qualification,
compile-time-resolvable text, and device address space. OpenCL-style
\texttt{printf} constrains its format to constant storage and its
\texttt{\%s} data to literal strings \cite{opencl-c}. A C++ cast neither
creates such storage nor makes runtime-generated bytes a literal.

\begin{lstlisting}
// Different capability tests, not interchangeable examples:
printf("%s", "literal");  // literal data
char data[8] = {'o', 'k', '\0'};
printf("%s", data);       // runtime/private string data
// A runtime-built first argument would test another capability.
\end{lstlisting}

Use a device-written value to make the second case truly runtime-dependent
in a capability test; otherwise optimization could fold it back into a
literal. Test a literal format plus private data, global data, a literal
accessed through a pointer, and a runtime format independently. Preserve
LLVM IR and final SPIR-V where available. LLVM address-space numbers and
SPIR-V storage-class enumerants are different naming systems and must not
be conflated in issue reports.

The existing extension
\path{SPV_EXT_relaxed_printf_string_address_space} is listed by LLVM as
relaxing format-string address spaces \cite{llvm-spirv}. This does not, by
itself, prove support for runtime \texttt{\%s} contents throughout a compiler,
validator, driver, and runtime stack. First identify which layer rejects the
required byte-output use case; then file a minimal implementation issue or a
precisely scoped extension request. This draft does not claim a new issue,
PR, or Khronos proposal has already been submitted.

A relevant upstream example is Intel graphics compiler issue 391, reporting
a chipStar-related \texttt{\%s} address-space failure; the retrieved page is
marked closed \cite{igc391}. It is related evidence, not proof that \ffmt's
reported problem is the same bug or is fixed in a particular release.

The project's \texttt{-O0} exclusions should likewise remain configuration-
specific observations \cite{ffmt-makefile}. The header names a SYCL
address-space transformation, but the assertion ``the pass does not run at
\texttt{-O0}'' needs pass-pipeline or IR evidence \cite{ffmt-source}.
A failure could instead involve lost provenance, a transformation defect,
argument lowering, or the runtime. Compare \texttt{-O0} and \texttt{-O2}
using the same minimized program before assigning a cause.

\section{Floating-point conversion and fast-math}

\subsection{Shortest conversion is only one component}

Dragonbox provides binary-to-decimal conversion for shortest round-trippable
representations \cite{dragonbox}. A formatter must additionally implement
notation selection, requested precision, rounding, signs, special values,
padding, and bounded writes. Shortest conversion does not prove correctness
of every fixed-precision path. A round-trip property is appropriate for
shortest finite conversion, not a universal test for deliberately rounded
output.

The numerical corpus should contain binary32 and binary64, both signed zeros,
subnormals, normal/subnormal boundaries, infinities, quiet and signaling NaN
encodings, extreme finite values, halfway cases, powers of ten, and adjacent
representable values. Store inputs as integer bit patterns rather than
reconstructing them through decimal parsing or arithmetic in the test driver.
State the expected handling of NaN spelling and sign; do not demand NaN
payload preservation from a textual format unless it is part of the contract.

\subsection{Robust classification has a limited, explicit scope}

Clang documents that fast-math can assume no NaNs or infinities and can ignore
signed zero \cite{clang-fp}. Therefore \texttt{x != x}, ordinary comparisons,
or a library classification call compiled under those assumptions are not a
sound universal diagnostic strategy. A suitable representation-level design
classifies integer encodings before floating arithmetic, subject to explicit
IEEE binary32/binary64 layout checks.

For example, for an IEEE binary32 encoding \texttt{bits}:
\begin{lstlisting}
constexpr bool finite_bits32(std::uint32_t bits) noexcept {
  return (bits & 0x7f800000u) != 0x7f800000u;
}
constexpr bool nan_bits32(std::uint32_t bits) noexcept {
  return (bits & 0x7f800000u) == 0x7f800000u &&
         (bits & 0x007fffffu) != 0;
}
\end{lstlisting}

Use \texttt{std::bit\_cast} or a supported equivalent to obtain an encoding
without violating aliasing rules; extract the sign from its sign bit.
For binary64, use the corresponding 11-bit exponent and 52-bit fraction.
Assert representation assumptions, and keep classification separate from
float-to-double conversion, which may have its own device capability and
special-value effects. The inspected \ffmt\ conversion code already uses
bit-level exponent/mantissa checks in relevant paths \cite{ffmt-source};
the work is to audit all paths, not to claim classification is wholly absent.

Integer classification cannot recover a subnormal already flushed to zero,
a sign discarded by earlier arithmetic, or a NaN changed before it reaches
the formatter. Nor does it make the rest of decimal conversion safe under
all fast-math transformations. Validate strict-math builds first; then test
fast-math, finite-only assumptions, signed-zero treatment, and device
flush-to-zero settings as distinct configurations. Mixed translation units,
link-time optimization, and runtime floating-point state belong in this
qualification.

The accompanying independent helper probe checks 20 binary encodings under
GCC 14.2.0 and Clang 17.0.0, each at \texttt{-O2} with and without
\texttt{-ffast-math}. It validates that small integer classifier only, not
\ffmt\ arithmetic, device lowering, or end-to-end fast-math correctness.

\section{Storage, limits, and generated code}

\subsection{Do not conflate format capacity and output capacity}

The inspected implementation has distinct limits: a default output capacity
of 128 characters, 256 bytes of format-text storage including the terminator,
and 16 placeholder entries in its buffer-path format object. The output size
is configurable with \texttt{FFMT\_BUFFER\_SIZE}; the buffer has 32 additional
scratch bytes \cite{ffmt-source}. These are implementation observations, not
a proof that every combination is safe. The format-text limit is therefore
at most 255 non-null characters for that object, not a 256-character output
promise.

\begin{lstlisting}
#define FFMT_BUFFER_SIZE 512
#include <ffmt/base.hpp>
\end{lstlisting}

Keep configuration macros consistent across translation units that share
header-defined types and functions; prefer target-wide build settings over
ad hoc per-file changes. Document whether each limit applies before or after
custom-formatter expansion. Width/precision representation, output capacity,
format length, and placeholder count require independent boundary tests.

Silent truncation is described by the README \cite{ffmt-readme}. A better
bounded-output contract would report bytes written and truncation, and
optionally required size when computing it is affordable. It should specify
terminator reservation, newline handling, embedded nulls, encoding, and
whether string width counts bytes or display columns. A \texttt{\%s}/
\texttt{fputs} sink cannot transparently transport embedded nulls. A byte
formatter is not automatically a Unicode-width-compatible standard formatter.

The previous source review identified writes that occur before length
clamping, notably in raw integer conversion \cite{ffmt-source}. An extra
scratch allowance is safe only when every helper's maximum write has been
proved to fit. Regressions should include a full logical buffer followed by
64-bit binary output, repeated fields after exhaustion, and very large fixed
precision. Retain sanitizer results before strengthening memory-safety claims.

\subsection{Private storage is not shared local memory}

An automatic per-invocation buffer is private logical storage. Depending on
optimization and indexing, it can become registers or per-thread local/scratch
memory; it is not automatically shared local memory (SLM), CUDA shared memory,
or AMD LDS. These resource categories must be reported separately. A larger
buffer can change stack or scratch traffic, spills, register pressure, and
occupancy, but none should be inferred numerically from the source declaration
alone \cite{cuda-memory,amdgpu-usage}.

Two sequential print calls need not allocate two simultaneously live buffers:
a compiler may reuse storage after the first lifetime ends. It is not required
to do so. Inlining, nested custom formatting, and sink lifetime requirements
can change the result. Compare one call, two sequential calls, and nested
formatting in the generated code. A sink must consume or copy temporary data
before its lifetime ends; host visibility of the output may occur later.

The same reasoning applies on CPUs: measure optimized stack frames and
allocation behavior rather than assuming a device-specific storage rule.
Compile-time parsing also does not guarantee zero runtime storage for its
parsed descriptor. Examine what is actually materialized at each optimization
level.

\subsection{Lookup tables and duplication}

Measure integer-only kernels, inclusion without use, shortest-float kernels,
and explicit-precision kernels separately. Inspect both pre-link objects and
the final device image for conversion tables and duplicate specializations.
Repeat with several call sites, multiple translation units, and link-time
optimization enabled and disabled. Tables may be shared, discarded, duplicated
across separate device images, or folded into other code; header source alone
cannot decide which outcome occurred. Record constant/read-only data separately
from code size and per-thread storage. Do not publish a table-size comment as
a measured final-binary cost.

\section{Validation and reproducibility}

\subsection{Run status is not pass status}

Table~\ref{tab:matrix} preserves the author's testing experience and adds
missing targets without manufacturing results. ``Reported run'' is not an
exact-output pass, and ``not independently run'' is not a zero-failure result.
The feasible compiler--runtime--device combinations matter, not a Cartesian
product of every programming model and vendor.

\begin{table}[!htbp]
\centering\small
\begin{tabularx}{\linewidth}{@{}P{2.35cm} Y P{2.25cm} Y@{}}
\toprule
Route & Configuration / device & Run evidence & Pass / issue evidence needed \\
\midrule
Host C++ & Linux and macOS CPUs & Author-reported & Exact compiler, libc++, libstdc++, ABI; both paths \\
OpenMP target & Clang / \texttt{icpx}; Intel and NVIDIA & Author-reported & Exact pairings; \texttt{-O0}/\texttt{-O2} exclusions and logs \\
SYCL & DPC++ on Intel & Author-reported & Device, runtime, driver, string cases \\
SYCL & AdaptiveCpp, \texttt{generic} & Author-reported & Actual backend/device and JIT details \\
CUDA C++ & Clang on NVIDIA & README-reported & Reproduce; separate \texttt{nvcc}/\texttt{nvc++} failures \\
HIP C++ & Native AMD route & Pending & Detection, annotations, sink, full test run \\
Kokkos & Each enabled execution space & Pending & Framework version and backend-specific runs \\
Arm host + GPU & Grace Hopper configuration & Pending & CPU run and GPU run recorded separately \\
chipStar & HIP/CUDA through SPIR-V & Pending & Distinct transport and address-space tests \\
\bottomrule
\end{tabularx}
\caption{Qualification ledger, not measured pass counts. No row was
independently executed with \ffmt\ for this revision. The CUDA entry and
compiler exclusions are repository descriptions \cite{ffmt-readme}.}
\label{tab:matrix}
\end{table}

AdaptiveCpp's \texttt{--acpp-targets=generic} names its generic compilation
flow, not a device or vendor \cite{acpp}. Kokkos similarly requires an
execution-space label. Retain the runtime-selected device, driver, architecture,
FP64 capability, host architecture, compiler and reference-library versions,
commit, flags, configuration macros, and effective output-buffer capacities.

For each configuration record build status, whether execution occurred,
positive-test pass/fail counts, negative-test acceptance/rejection counts,
skips, expected failures, unexpected passes, and linked bug identifiers.
Use ``not measured'' rather than zero where no log exists. Unique logical
cases and total executions across configurations are different quantities.

\subsection{Positive, negative, and runtime-invalid tests}

Negative compilation tests need at least two distinct suites. Invalid C++
format strings should be rejected by the selected standard reference and
\ffmt. Valid formats outside a backend's supported subset should be accepted
by the reference but deliberately rejected by that backend. Combining these
suites would incorrectly treat a backend limitation as a language error.

\begin{table}[!htbp]
\centering\small
\begin{tabularx}{\linewidth}{@{}P{3.0cm} Y Y@{}}
\toprule
Class & Representative cases & Expected evidence \\
\midrule
Invalid static syntax/type & Unmatched braces; mixed auto/manual indices; missing argument; integer with \texttt{\{:s\}} & Intended compile diagnostic, not merely any nonzero compiler exit \\
Valid but unsupported & Binary, custom fill, or dynamic precision on a restricted path & Explicit backend rejection; standard reference acceptance \\
Runtime-invalid values & Negative dynamic precision, excessive runtime width, invalid string lifetime & Documented runtime policy; cannot universally reject at compile time \\
Boundary and delivery & Full buffer; terminator/newline; percent expansion; concurrent records & Sanitizer/bounds evidence plus exact output or classified transport failure \\
\bottomrule
\end{tabularx}
\caption{Separate validation obligations. Examples are proposed tests, not
newly measured outcomes.}
\label{tab:negative}
\end{table}

A negative-test harness must retain diagnostics, first compile a positive
control, and distinguish a deliberate format rejection from a missing header,
compiler crash, unsupported target, or unrelated error. The inspected Makefile
currently accepts a nonzero compiler exit in its negative-test loop, so
reason-aware validation is a necessary improvement \cite{ffmt-makefile}.
Do not reject extra unused arguments merely by counting placeholders; consult
the selected format specification and test the actual intended rule.

For supported formatting, compare exact bytes under a controlled locale and
reference version. Include integer limits, signs, alternate forms, escaped
braces, empty strings, positional reuse, and custom-formatter expansion.
Combine bounded-write sanitizers with differential output tests; matching
text alone does not establish safe intermediate writes. Fuzz format grammar
and values, retaining seeds and minimized reproducers.

\subsection{Device-output and resource tests}

First isolate formatting by writing results to a test sink or device memory
and reading them back. Then test the real output primitive using unique
record IDs, counts, and delimiters. Distinguish permitted reordering from
interleaving, duplication, truncation, and missing records; sorting output
must not hide those errors. One call per record is a useful design choice,
not a universal atomicity guarantee \cite{opencl-c}.

CUDA documents a circular output buffer, delayed flushing, and host-side
final formatting \cite{cuda-printf}. Therefore a successful kernel completion
or print return value is not, alone, evidence that every expected byte reached
the captured stream. Synchronize and drain according to the runtime, retain
launch/asynchronous errors, and test output-volume limits separately from
\ffmt's local capacity. Repeat address-space and argument-ABI cases at
\texttt{-O0} and optimized levels before relying on compiler workarounds.

\subsection{Build, sanitizers, and coverage}

The Makefile supplies host, negative, coverage, and accelerator test routes,
including optimization-specific exclusions \cite{ffmt-makefile}. The following
are proposed commands for a complete checkout, not a record of executed
\ffmt\ tests:

\begin{lstlisting}[language=bash]
git rev-parse HEAD
c++ --version
make clean
make test-host CXX=g++ HOST_OPT=-O2
make test-negative CXX=g++
make clean
make test-host CXX=g++ HOST_OPT=-O1 \
  EXTRA_CXXFLAGS="-g -fno-omit-frame-pointer -fsanitize=address,undefined"
\end{lstlisting}

Host sanitizers do not qualify GPU address spaces or runtimes. Retain device
checks separately. Report line and branch coverage for each instrumented
path, excluded code and compiler workarounds, and the exact test denominator;
coverage percentages are not semantic coverage. In particular, switching to
\texttt{-O2} does not erase the Makefile's optimized-build exclusions.

A CMake/CTest front end is proposed engineering work, not a claimed repository
feature: define an interface library, separate build directories/toolchain
presets, explicit emission options, compiled README examples, and diagnostic-
aware negative tests. Remove developer-local compiler paths and preserve the
existing Make targets until both routes exercise the same cases. Audit
upstream attribution and licensing when packaging the port.

\section{Performance evaluation plan}

The research question is not simply whether a print call is ``fast.'' Measure
three layers independently, with a no-print kernel as an additional baseline.
No timing or resource result is asserted in this draft.

\begin{table}[!htbp]
\centering\small
\begin{tabularx}{\linewidth}{@{}P{2.6cm} Y Y@{}}
\toprule
Experiment & Fair comparison & Measurements \\
\midrule
Compilation & Same corpus; ordinary \fmtlib, \texttt{FMT\_COMPILE}, standard facilities, and \ffmt\ configurations labeled separately & Front-end and device compilation, link time, peak memory, cold JIT versus warm execution \\
Formatting only & Same values, semantics, capacity, and observable memory sink; no terminal I/O & Time per value/byte, output checksum, code and read-only data size \\
End-to-end output & Same kernel work, record count, output volume, capture, synchronization, and drain policy & Kernel duration and wall time; lost records; runtime-buffer capacity \\
Resource scaling & One/two calls; integer/float; several buffer sizes and translation units & Registers, spills, stack/scratch, occupancy, SLM/shared memory, final-image duplication \\
\bottomrule
\end{tabularx}
\caption{Evaluation design. Results remain to be collected.}
\label{tab:performance}
\end{table}

On CPUs, compare against \texttt{fmt::format\_to} and
\texttt{std::format\_to} into equivalent preallocated storage, with
\texttt{snprintf} on the semantically common subset. Do not confuse avoiding
\texttt{std::format}'s returned-string allocation with a faster conversion
algorithm. An ordinary \fmtlib\ build and a \texttt{FMT\_COMPILE} build are
different baselines \cite{fmt-api}.

On devices, benchmark raw backend \texttt{printf} against \ffmt\ only where
the outputs are intended to match. A shortest conversion and a six-digit
\texttt{\%g} conversion are different tasks. A host-only library is not a
GPU benchmark result. Use runtime-generated inputs and an observable sink to
prevent constant folding or dead-code elimination. Capture warm-up policy,
repetitions, medians and spread, compiler flags, and input distributions.

For generated-code analysis retain LLVM IR or SPIR-V, PTX and final machine
code where available, plus compiler resource reports. Separate static metadata
from measured peak live storage. Evaluate code-size growth with the number of
distinct format strings, argument types, and custom expansions; compile-time
specialization may trade execution work for compilation cost and binary size.

\section{Extensions and remaining limitations}

\subsection{Framework formatters and output adapters}

The inspected header already defines formatters for
\texttt{sycl::range}, \texttt{sycl::id}, \texttt{sycl::item}, and
\texttt{sycl::nd\_item} \cite{ffmt-source}. Future work should extend and
validate that facility rather than present all SYCL pretty-printing as new.
Document custom nesting, positional-index restrictions, lifetime rules, and
how a formatter inherits the selected backend's unsupported features.

Kokkos integration should start with a qualified emission adapter using
\texttt{Kokkos::printf}, which is documented for host and device use, and then
test each enabled execution space \cite{kokkos-printf}. A wrapper cannot
remove the underlying runtime's string restrictions or repair missing
host/device annotations transitively. For \texttt{Kokkos::View}-like objects,
prefer bounded metadata summaries such as rank, extents, strides, and memory
space. Printing contents needs an explicit policy for accessibility, maximum
length, synchronization, and whether host/device copies occur.

\subsection{Python, Fortran, and a stronger bounded-output core}

Python is useful first as host-side test orchestration: generate typed C++
cases, preserve input bit patterns, capture manifests and logs, and summarize
qualification results. Python formatting is not a silent replacement for the
chosen C++ oracle. A host binding to a formatting core is a separate project
from formatting inside a Python-launched device kernel.

Fortran interoperation would need a deliberately small interoperable boundary,
for example a typed wrapper around caller-owned byte storage with explicit
length and status. C++ template packs, C varargs, and Fortran character
arguments are not interchangeable. Device-language integration would require
its own compiler and offload tests; this draft claims no existing binding.

A future core API could separate bounded formatting from emission entirely,
with a size-aware byte sink and explicit errors. Such an interface would
simplify memory-safety testing, embedded-null handling, and formatting-only
benchmarks. It does not remove backend transport restrictions. Locales,
Unicode display width, wide strings, extended floating types, chrono/ranges,
and unrestricted standard formatter customization should remain outside any
unqualified full-compatibility claim until individually specified and tested.

\section{Conclusion}

\ffmt's useful abstraction is a common validated interface over two different
output capabilities. Producing bytes in the library gives control over their
representation; delegating values to \texttt{printf} trades some of that
control for compatibility with a restricted sink. The distinction explains
both the architecture and its test obligations. A publishable evaluation
should connect each supported feature to versioned byte-level evidence and
each backend to build, execution, and output-delivery logs. That evidence,
rather than a single undifferentiated support checkmark, is the basis for
claims of portability, safety, and performance.

\section*{Availability and publication status}

Project: \url{https://github.com/TApplencourt/ffmt}.
The implementation credits \fmtlib\ and Dragonbox \cite{ffmt-readme,dragonbox}.
This is an expanded technical working draft, not an anonymized or
conference-formatted submission. The companion revision notes contain the
ISC 2027 requirements, team work items, evidence gaps, and source-review
ledger. Replace mutable source references with an immutable commit and
reproducible artifact before publication. Follow venue instructions for
self-citations, identifying repository links, and substantive AI assistance.

\begin{thebibliography}{99}
\footnotesize
\setlength{\itemsep}{0.15em}
\setlength{\parskip}{0pt}

\bibitem{ffmt-readme}
\ffmt\ project.
\emph{Repository README and backend descriptions}, main branch.
Accessed 21 October 2026.
\url{https://github.com/TApplencourt/ffmt}.

\bibitem{ffmt-source}
\ffmt\ project.
\emph{include/ffmt/base.hpp}, main branch.
Accessed 21 October 2026; mutable source, not a pinned release.
\url{https://raw.githubusercontent.com/TApplencourt/ffmt/main/include/ffmt/base.hpp}.

\bibitem{ffmt-makefile}
\ffmt\ project.
\emph{Makefile}, main branch. Accessed 21 October 2026.
\url{https://raw.githubusercontent.com/TApplencourt/ffmt/main/Makefile}.

\bibitem{format-intro}
Mark de Wever. ISO C++ committee paper P2930R0.
\emph{Formatter specializations for the standard library}, 2023.
\url{https://www.open-std.org/jtc1/sc22/wg21/docs/papers/2023/p2930r0.html}.

\bibitem{print-intro}
ISO C++ committee paper P3107R5.
\emph{Permit an efficient implementation of std::print}, 2024.
\url{https://open-std.org/jtc1/sc22/wg21/docs/papers/2024/p3107r5.html}.

\bibitem{compliance}
C++ working draft.
\emph{Implementation compliance [compliance]}.
Live draft accessed 21 October 2026.
\url{https://eel.is/c++draft/compliance}.

\bibitem{fmt-api}
\fmtlib\ project.
\emph{API documentation}, including configuration and compatibility notes.
Accessed 21 October 2026.
\url{https://fmt.dev/latest/api/}.

\bibitem{format-functions}
ISO C++ working draft N4950.
\emph{Formatting functions [format.functions]}, 2023.
HTML rendering of the committee draft.
\url{https://timsong-cpp.github.io/cppwp/n4950/format.functions}.

\bibitem{n4950-format}
ISO C++ working draft N4950.
\emph{Standard format specifiers [format.string.std]}, 2023.
\url{https://timsong-cpp.github.io/cppwp/n4950/format.string.std}.

\bibitem{n4950-charconv}
ISO C++ working draft N4950.
\emph{Primitive numeric output conversion [charconv.to.chars]}, 2023.
\url{https://timsong-cpp.github.io/cppwp/n4950/charconv.to.chars}.

\bibitem{current-charconv}
C++ working draft.
\emph{Primitive numeric output conversion [charconv.to.chars]}.
Live draft accessed 21 October 2026.
\url{https://eel.is/c++draft/charconv.to.chars}.

\bibitem{p3505}
Victor Zverovich and Junekey Jeon.
\emph{Fix the default floating-point representation in std::format}.
ISO C++ committee paper P3505R3, 2025.
\url{https://open-std.org/JTC1/SC22/WG21/docs/papers/2025/p3505r3.html}.

\bibitem{dragonbox}
Junekey Jeon.
\emph{Dragonbox: reference implementation and algorithm description}.
Accessed 21 October 2026.
\url{https://github.com/jk-jeon/dragonbox/blob/master/README.md}.

\bibitem{opencl-c}
Khronos Group.
\emph{The OpenCL C Specification}, section 6.15.14 (printf).
Unified specification v3.1.2, 17 September 2026. Accessed 21 October 2026.
\url{https://registry.khronos.org/OpenCL/specs/unified/html/OpenCL_C.html#printf}.

\bibitem{llvm-spirv}
LLVM project.
\emph{User Guide for SPIR-V Target}, supported extensions.
Accessed 21 October 2026.
\url{https://llvm.org/docs/SPIRVUsage.html}.


\bibitem{acpp}
AdaptiveCpp project.
\emph{Compilation flows: generic SSCP compilation}.
Accessed 21 October 2026.
\url{https://github.com/AdaptiveCpp/AdaptiveCpp/blob/develop/doc/compilation.md}.

\bibitem{hip}
AMD.
\emph{HIP C++ language extensions}.
Accessed 21 October 2026.
\url{https://rocm.docs.amd.com/projects/HIP/en/latest/how-to/hip_cpp_language_extensions.html}.


\bibitem{format-string}
C++ working draft. \emph{Class template basic\_format\_string [format.fmt.string]}.
Live draft accessed 21 October 2026.
\url{https://eel.is/c++draft/format.fmt.string}.

\bibitem{intel-printf}
Intel. \emph{SYCL Breaking Changes in Intel oneAPI DPC++/C++ Compiler 2026},
section on removal of the variadic printf implementation. Accessed 21 October 2026.
\url{https://www.intel.com/content/www/us/en/developer/articles/technical/sycl-breaking-changes-oneapi-dpcpp-compiler-2026.html}.

\bibitem{cuda-printf}
NVIDIA. \emph{CUDA Programming Guide: C++ Language Support}, formatted output.
Accessed 21 October 2026. Distinguish SIMT from tile-code restrictions.
\url{https://docs.nvidia.com/cuda/cuda-programming-guide/05-appendices/cpp-language-support.html}.

\bibitem{amdgpu-printf}
LLVM project. \emph{AMDGPUEmitPrintf.cpp}, printf lowering to device-library calls.
Main branch, accessed 21 October 2026.
\url{https://raw.githubusercontent.com/llvm/llvm-project/main/llvm/lib/Transforms/Utils/AMDGPUEmitPrintf.cpp}.

\bibitem{kokkos-printf}
Kokkos project. \emph{Kokkos::printf API documentation}.
Accessed 21 October 2026.
\url{https://kokkos.org/kokkos-core-wiki/API/core/utilities/printf.html}.

\bibitem{chipstar}
CHIP-SPV project. \emph{chipStar repository and project description}.
Accessed 21 October 2026.
\url{https://github.com/CHIP-SPV/chipStar}.

\bibitem{grace-hopper}
NVIDIA. \emph{NVIDIA Grace Hopper Superchip Architecture In-Depth}.
Accessed 21 October 2026.
\url{https://developer.nvidia.com/blog/nvidia-grace-hopper-superchip-architecture-in-depth}.

\bibitem{igc391}
Intel graphics compiler issue 391.
\emph{printf crashes with \%s argument in UniformConstant storage class (SPIR-V)}.
Opened 10 March 2026; retrieved page marked closed. Accessed 21 October 2026.
\url{https://github.com/intel/intel-graphics-compiler/issues/391}.

\bibitem{clang-fp}
LLVM project. \emph{Clang Compiler User's Manual}, floating-point options.
Accessed 21 October 2026.
\url{https://clang.llvm.org/docs/UsersManual.html}.

\bibitem{cuda-memory}
NVIDIA. \emph{CUDA Programming Guide: Writing SIMT Kernels}, memory hierarchy.
Accessed 21 October 2026.
\url{https://docs.nvidia.com/cuda/cuda-programming-guide/02-basics/writing-cuda-kernels.html}.

\bibitem{amdgpu-usage}
LLVM project. \emph{User Guide for AMDGPU Backend}, address spaces and scratch.
Accessed 21 October 2026.
\url{https://llvm.org/docs/AMDGPUUsage.html}.

\end{thebibliography}
\end{document}
